\documentclass[11pt]{article}
\usepackage[a4paper,margin=28mm]{geometry}
\usepackage{amsmath,amssymb,amsthm,mathtools,hyperref}
\newcommand{\ddc}{dd^c}
\newcommand{\PSH}{\operatorname{PSH}}
\newcommand{\MA}{\operatorname{MA}}
\newcommand{\PP}{\mathbb P}
\newcommand{\BB}{\mathcal B}
\newtheorem{theorem}{Theorem}
\newtheorem{lemma}[theorem]{Lemma}
\newtheorem{proposition}[theorem]{Proposition}
\theoremstyle{remark}\newtheorem{remark}[theorem]{Remark}
\title{Perron interpolation and metric actions on Hopf manifolds}
\author{}
\date{8 October 2026}
\begin{document}
\maketitle
\noindent\textbf{Checked restricted theorem.} Complete proof review,
independent statement-only reconstruction and final comparison establish
the A--H results in this article. The original review and the first
final review remain INCONCLUSIVE in the retained records; a fresh final
review supplies valid bindings for the precisely scoped imported
foundations. The general singularity class programme remains open.
\begin{abstract}
On the diagonal Hopf manifold of every dimension at least three, we study
canonical potentials invariant under the elliptic fiber action. The full
Hermitian Perron problem for two such endpoints descends to the projective
base. Its interpolation is the unique geodesic of a complete metric on the
invariant finite energy potentials. The associated metric action agrees
with the raw Hermitian integral action on smooth invariant paths. Smooth
Perron paths are stationary under every smooth variation, by compact group
averaging. A diagonal projective family gives explicit nonconstant examples
and exact distance and action formulas. The metric acts on potentials;
their ordinary model projection in this finite energy domain is zero.
An explicit unbounded family has finite squared energy exactly below a
computed exponent threshold and contains distinct ordinary types.
Constant Perron translations are nevertheless saddles for the full raw
action, with infinitely many negative directions and strictly smaller
lengths among nonbasic paths.
\end{abstract}

We use $d^c=\frac{i}{4\pi}(\bar\partial-\partial)$ and
$\ddc=\frac{i}{2\pi}\partial\bar\partial$, ordinary wedge powers,
and genuine canonical quasi-psh representatives, including their poles.
For bounded potentials we use the local products with smooth corrections
\cite[Section 2.1 and Proposition 2.1]{BGL}, including their plurifine
locality \cite{BTfine}. The unbounded basic products used here are proved
from canonical finite-level identities below; no BMP-dependent global
nonclosed product theorem is applied on $X$. Constants are retained in
the potential space and have positive distance.

\section{The elliptic bundle and its basic potentials}
Fix $n\geq3$, put $d=n-1$, and consider
\[
 X=(\mathbb C^n\setminus\{0\})/\langle z\mapsto ez\rangle,
 \qquad \pi:X\longrightarrow Y=\PP^d.
\]
The free deck action is properly discontinuous and has the compact radial
fundamental strip $1\leq\lVert z\rVert\leq e$. Thus $X$ is connected and
compact. Its fibers are the elliptic group
$E=\mathbb C/(\mathbb Z+2\pi i\mathbb Z)$, acting transitively by
scalar multiplication. Topologically $X=S^{2n-1}\times S^1$; its first
Betti number is one, whereas that of $\PP^d\times E$ is two.
Its real cohomology in every intermediate positive even degree vanishes.
A product of two positive-dimensional compact complex manifolds would
have the nonzero pullback of either factor's top-degree class in such
a degree, since projection to that factor has a smooth section.
Thus $X$ admits no such product decomposition.

On the cover write
\[
 t=\log\lVert z\rVert^2,\qquad
 \alpha=\ddc t=\pi^*\kappa,\qquad
 \beta=dt\wedge d^ct,\qquad h=\alpha+\beta,
\]
where $\int_Y\kappa^d=1$. The deck map adds two to $t$, so these
forms descend. The form $\alpha$ is positive in base directions, and
$\beta$ is semipositive of rank one and positive vertically. Hence $h$
is Hermitian. Moreover
\begin{equation}\label{eq:geometry}
 dh=-dt\wedge\alpha=-dt\wedge h\ne0,
 \qquad \int_{\pi^{-1}(y)}\beta=2,
 \qquad \int_Xh^n=2n.
\end{equation}
For the middle equality take a fiber coordinate $\zeta=a+ib$. Then
$t=2a+\text{constant}$ and $\beta=da\wedge db/\pi$, with periods
$1$ and $2\pi$. The last equality follows from $\beta^2=0$ and
$\alpha^{d+1}=0$. Both $h$ and these normalizations are invariant under
$E$. An invariant genuine potential is called basic and is exactly $\pi^*v$.

\begin{lemma}\label{lem:descent}
For genuine potentials on $Y$, and for their complexifications on
$Y\times S$, where $S$ is a strip, positivity satisfies
\[
 \operatorname{pr}_X^*h+\ddc(\pi\times\mathrm{id})^*V\geq0
 \quad\Longleftrightarrow\quad
 \operatorname{pr}_Y^*\kappa+\ddc V\geq0.
\]
For every genuine $v\in\PSH(Y,\kappa)$ the actual measure is
\begin{equation}\label{eq:measure}
 \MA_h^{\mathrm{ext}}(\pi^*v)
     =n\beta\wedge\pi^*\MA_\kappa(v).
\end{equation}
Consequently, for every nonnegative Borel $F$,
\begin{equation}\label{eq:fiber-integration}
 \int_X\pi^*F\,\MA_h^{\mathrm{ext}}(\pi^*v)
       =2n\int_YF\,\MA_\kappa(v).
\end{equation}
\end{lemma}
\begin{proof}
Locally $\beta=\frac{i}{2\pi}\partial t\wedge\bar\partial t$ has
a nonzero vertical coefficient. Every vector on $Y\times S$ has a unique
lift annihilated by $\partial t$. Testing the positive current on those
smooth lifts removes $\beta$ and gives positivity of the base current.
Its distributional coefficients are independent of the fiber coordinate;
testing with a nonnegative compact fiber density proves the assertion
on the base. The reverse implication adds $\beta\geq0$. Local
holomorphic sections give the canonical quasi-psh representative of a
descended invariant function, so this argument includes singular potentials.

For bounded $v$, local Bedford--Taylor multiplication and the smooth
correction polynomial give \eqref{eq:measure}: the only nonzero top term
contains one $\beta$ and $d$ base factors. For general $v$, apply that
identity to $v_j=\max(v,-j)$ and multiply by
$\mathbf1_{\{v>-j\}}$. The increasing finite-level measures on both
sides give \eqref{eq:measure} without a product on the pole set.
Fiber integration of the smooth vertical factor, of mass two, proves
\eqref{eq:fiber-integration}.
\end{proof}

\section{The full Perron envelope and its completion}
The Hermitian Perron segment uses all strip subgeodesics, with joint
boundary upper bounds given by the endpoints. Basicness is not an
assumption on its competitors.
\begin{proposition}\label{prop:perron}
For two basic endpoints, whenever the base Perron family is nonempty,
the full Hermitian Perron envelope is the pullback of the base envelope.
For $v_0,v_1\in\mathcal E^2(Y,\kappa)$ it attains both endpoints
genuinely in $L^1(X)$.
\end{proposition}
\begin{proof}
The full competitor family is stable under the compact fiber group.
Its supremum and usc regularization are therefore invariant. The locally
bounded upper chord bound and the existence of a base competitor make
the resulting envelope a genuine basic potential. Lemma~\ref{lem:descent}
makes its descent a base subgeodesic with the same boundary bounds.
Conversely, every base competitor pulls back to a full competitor.
This proves equality of the two envelopes; usc regularization commutes
with pullback in submersion charts.

For smooth endpoints a competitor is
$\max\{v_0-Ct,v_1-C(1-t)\}$, with
$C\geq\sup_Y|v_0-v_1|$. For arbitrary $\mathcal E^2$ endpoints,
the base Perron construction is nontrivial and has capacity endpoint
attainment \cite[Theorem 5.2 and Corollary 5.4]{Darvas}. It therefore
attains them in base $L^1$. Equation~\eqref{eq:geometry} gives
\[
 \int_X|\pi^*(v-w)|\frac{h^n}{2n}
       =\int_Y|v-w|\kappa^d,
\]
which proves the claimed endpoint attainment on $X$. Every prescribed
smooth positive probability $dV$ has a bounded positive density ratio
with $h^n/(2n)$ on compact $X$, so this convergence also holds in the
original task's $L^1(X,dV)$, with the same canonical endpoints.
\end{proof}

Let $\BB_\infty$ be the smooth basic potentials with strictly positive
$h+\ddc u$, and set
\[
 \BB_2=\left\{\pi^*v:\ v\in\PSH(Y,\kappa),\quad
 \int_Y\MA_\kappa(v)=1,\quad
 \int_Y|v|^2\MA_\kappa(v)<\infty\right\}.
\]
By \eqref{eq:fiber-integration}, these are precisely the basic potentials
of actual mass $2n$ and finite squared energy on $X$.
Define the distance on $\BB_\infty$ by the infimum of raw Hermitian
lengths of smooth \emph{basic} paths. The raw metric restricted to those
tangents is exactly
\begin{equation}\label{eq:metric}
 G^X_{\pi^*v}(\pi^*\xi,\pi^*\eta)
       =2n\int_Y\xi\eta(\kappa+\ddc v)^d.
\end{equation}

\begin{theorem}\label{th:completion}
The metric completion of $\BB_\infty$ is $\BB_2$, with distance
\begin{equation}\label{eq:distance}
 d_B(\pi^*v,\pi^*w)
   =2\pi\sqrt{2n}\,\widetilde d_D
                 \left(\frac{v}{2\pi},\frac{w}{2\pi}\right).
\end{equation}
Here $\widetilde d_D$ is Darvas' completed metric for the background
$\kappa$ and operator $i\partial\bar\partial$.
Its unique constant-speed geodesic is the full Hermitian Perron segment.
\end{theorem}
\begin{proof}
The operator conversion is
$\kappa+i\partial\bar\partial(v/(2\pi))=\kappa+\ddc v$.
Thus \eqref{eq:metric} gives \eqref{eq:distance} on smooth potentials.
Darvas' completion theorem applies on the connected compact K\"ahler
manifold $Y$ \cite[Theorems 6.1 and 9.2]{Darvas}. Constant positive
rescaling preserves its completeness and unique geodesics, giving the
completion on genuine $\mathcal E^2$ representatives. Proposition
\ref{prop:perron} identifies that geodesic with the full Hermitian
envelope. In particular
$d_B(\pi^*v,\pi^*(v+c))=\sqrt{2n}|c|$.
\end{proof}

For absolutely continuous curves in this completed metric set
\begin{equation}\label{eq:action}
 \mathcal A_B(\gamma)=\frac12\int_0^1|\dot\gamma|_{d_B}^2\,dt.
\end{equation}
Every connecting curve has action at least half the squared endpoint
distance, by the length inequality and Cauchy--Schwarz. The unique
constant-speed Perron segment attains equality and is the unique minimizer.

On a smooth basic path \eqref{eq:action} is the raw integral action.
For completeness, the metric derivative can be checked from the
smooth-endpoint formula \cite[(15)]{Darvas}. If $v_{s+\delta}-v_s
=\delta\dot v_s+O(\delta^2)$ in $C^2$, its short base Perron segment
lies between its chord and that chord minus
$C\delta^2\tau(1-\tau)$. The lower path is a subgeodesic by the
Schur inequality and uniform strict positivity near $v_s$.
Its initial velocity is consequently
$v_{s+\delta}-v_s+O(\delta^2)$ uniformly. Formula (15) gives the
metric derivative $\bigl(\int_Y\dot v_s^2\MA_\kappa(v_s)\bigr)^{1/2}$
on the base. Equation~\eqref{eq:metric} gives the stated raw speed on $X$.
For a path only $C^1$ in time the Taylor remainder is $o(\delta)$,
and the same argument has the same derivative limit.

The completed action is defined by metric speed. The initial-velocity
formula has a bounded-potential counterexample outside its smooth
hypotheses \cite[Section 7]{Darvas}; it is not used to define a
Riemannian metric at every point of $\BB_2$.

\section{Stationarity under every smooth variation}
\begin{theorem}\label{th:stationary}
Every smooth basic Perron path with strictly positive spatial forms is
stationary for the original raw Hermitian action under all smooth
variations with fixed endpoints.
\end{theorem}
\begin{proof}
Write $u_t=\pi^*v_t$ and let $L(w)$ be the first variation of
\[
 \mathcal S_X(u)=\frac12\int_0^1\int_X\dot u_t^2(h+\ddc u_t)^n\,dt.
\]
Uniform positivity on $[0,1]\times X$ makes this derivative well defined
for every smooth endpoint-zero variation. The action is invariant under
$E$ and the path is fixed by $E$. Hence $L(g^*w)=L(w)$ and
\[
 L(w)=L\left(\int_Eg^*w\,dg\right).
\]
The averaged variation is smooth basic, say $\pi^*q_t$.
Equation~\eqref{eq:fiber-integration}, applied before differentiating,
reduces its first variation to $2n$ times the closed-base one.
Put $\eta_t=\kappa+\ddc v_t$. Time integration by parts and closed
spatial integration by parts give
\[
 L_Y(q)=\int_0^1\int_Y q_t\left[
 -\ddot v_t\eta_t^d
 +d\,d\dot v_t\wedge d^c\dot v_t\wedge\eta_t^{d-1}\right]dt.
\]
Indeed $\ddc(f^2)=2f\ddc f+2df\wedge d^cf$ cancels the two
terms containing $\dot v_t\ddc\dot v_t$. The bracket vanishes by
the smooth base homogeneous Monge--Amp\`ere equation, or its Schur
complement. Thus $L_Y(q)=0$ and $L(w)=0$ for every variation.
\end{proof}

This theorem concerns stationarity, while the minimization in
\eqref{eq:action} is in $\BB_2$. The length infimum in
Theorem~\ref{th:completion} is over basic paths. A global raw-action
minimum among all nonbasic paths already fails for constant endpoints,
as the next theorem shows. The torsion is still present:
here $\ddc\beta=-\alpha^2$,
but
\[
 \ddc(h+\ddc\pi^*v)^{n-1}
   =-(n-1)\alpha^2\wedge\pi^*\eta^{n-2}=0
\]
by the dimension of the base. The group averaging proves the full
first-variation assertion without suppressing torsion in a general metric.

\section{Stationary translations are saddles}
The distinction between stationarity and minimization is substantive.
Write $r=\log\lVert z\rVert^2$ modulo two. Every smooth two-periodic
function of $r$ is a genuine global smooth function on $X$. For a path
$u(t,r)$ one has
\[
 h+\ddc u=(1+u_r)\alpha+(1+u_{rr})\beta .
\]
Thus both coefficients positive imply strict spatial positivity.
Let $\langle F\rangle=\frac12\int_0^2F(r)\,dr$.
The fiber integral of $\beta$ is two, so the exact raw action is
\begin{equation}\label{eq:radial-action}
 \mathcal S_X(u)=n\int_0^1
   \left\langle u_t^2(1+u_r)^d(1+u_{rr})\right\rangle dt.
\end{equation}
The average is independent of the choice of base section, which only
shifts the start of a full period.

\begin{theorem}\label{th:saddle}
For every $c\ne0$, the full Hermitian Perron segment $u_t=ct$ is
stationary under all smooth fixed-endpoint variations, but its raw-action
Hessian has an infinite-dimensional negative subspace.
The infimum of raw Hermitian lengths over all smooth positive paths
from $0$ to $c$ is strictly smaller than $d_B(0,c)=\sqrt{2n}|c|$.
\end{theorem}
\begin{proof}
Every strip subgeodesic is convex in real time at a fixed spatial point,
so its endpoint upper bounds imply $u_t\leq ct$. The path $ct$ itself
is a competitor. It is therefore the full Perron segment.
It is smooth basic and stationary by Theorem~\ref{th:stationary}.

Take a smooth periodic $w(t,r)$ with $w(0,r)=w(1,r)=0$.
The paths $ct+\varepsilon w$ are spatially strictly positive for small
$\varepsilon$. Expanding \eqref{eq:radial-action}, and integrating the
terms $w_rw_{rr}$ in $r$ and $w_tw_{rr}$ in both $r$ and $t$, gives
\[
 \left.\frac{d^2}{d\varepsilon^2}\mathcal S_X(ct+\varepsilon w)
 \right|_{\varepsilon=0}
 =2n\int_0^1\left\langle
       (w_t+dcw_r)^2-\frac{d(d+1)}2c^2w_r^2
                    \right\rangle dt.
\]
Indeed $\int\langle w_tw_{rr}\rangle dt
=-\frac12[\langle w_r^2\rangle]_0^1=0$.
For every positive integer $k$, set
\[
 w_k(t,r)=\sin(\pi t)\cos\bigl(\pi k(r-dct)\bigr).
\]
Then $w_{k,t}+dcw_{k,r}=\pi\cos(\pi t)\cos(\pi k(r-dct))$,
and the displayed quadratic form is exactly
\begin{equation}\label{eq:saddle-mode}
 \frac{n\pi^2}{2}\left(1-\frac{d(d+1)}2c^2k^2\right).
\end{equation}
It is negative for all sufficiently large $k$. Distinct positive
integer modes and their derivatives are orthogonal under $r$-averaging.
The quadratic form is therefore negative on every nonzero finite
combination of those large modes.
Every such combination has admissible sufficiently small variations.

Their action is strictly less than $\mathcal S_X(ct)=nc^2$.
On the unit time interval their raw length satisfies
$L_{\mathrm{raw}}^2\leq2\mathcal S_X$ by Cauchy--Schwarz.
Thus an admissible path has length strictly less than $\sqrt{2n}|c|$.
These raw-action competitors are not required to be strip subgeodesics.
No definiteness or completeness of the full length infimum is asserted.
\end{proof}

For a completely explicit threefold example take
\[
 u_\varepsilon(t,r)=t+\varepsilon\sin(\pi t)
                         \cos\bigl(\pi(r-2t)\bigr).
\]
Its spatial coefficients are at least $1-\pi|\varepsilon|$ and
$1-\pi^2|\varepsilon|$. Direct periodic averaging in
\eqref{eq:radial-action} gives the exact polynomial
\begin{equation}\label{eq:saddle-polynomial}
 \mathcal S_X(u_\varepsilon)
    =3-\frac{3\pi^2}{2}\varepsilon^2
                     +\frac{111\pi^4}{64}\varepsilon^4.
\end{equation}
To check the quartic coefficient, write $f=\sin(\pi t)$ and $p=f'$.
The averaged quadratic coefficient before multiplication by $n=3$ is
$(p^2-3\pi^2f^2-2\pi^2fp)/2$; the quartic one is
$\pi^2f^2(p^2+6\pi^2fp+12\pi^2f^2)/8$.
All odd coefficients vanish. Use
$\int fp=\int f^3p=0$, $\int f^2=1/2$, $\int p^2=\pi^2/2$,
$\int f^4=3/8$ and $\int f^2p^2=\pi^2/8$.
The choice $\varepsilon=1/100$ is strictly positive spatially and
has action $2.9985212487\ldots<3$.

\section{Explicit nonconstant paths}
Choose real $a_0,\ldots,a_{n-1}$, not all equal, and set
\[
 v_t([z])=\log\frac{\sum_{j=0}^{n-1}e^{ta_j}|z_j|^2}
                        {\sum_{j=0}^{n-1}|z_j|^2}.
\]
On the strip, the holomorphic map
$F([z],s)=[e^{a_0s/2}z_0:\cdots:e^{a_{n-1}s/2}z_{n-1}]$
satisfies
\[
 \operatorname{pr}_Y^*\kappa+\ddc V=F^*\kappa.
\]
The form has rank $d$, its spatial slices are strictly positive, and
its $(d+1)$st power vanishes. The bounded base Dirichlet uniqueness
\cite[Section 2.2, (16)--(17)]{Darvas} identifies this smooth path
with the Perron segment. Its pullback is the full Hermitian segment
and is stationary under every smooth variation by
Theorem~\ref{th:stationary}.

Under $\kappa^d$, the variables
$x_j=|z_j|^2/\sum|z_j|^2$ have the uniform simplex distribution.
To see this, take a standard complex Gaussian vector, divide it by its
norm, and use its unitary invariant projective distribution. Its squared
coordinates are independent exponential variables before division.
Thus $\int x_i^2\kappa^d=2/[n(n+1)]$ and
$\int x_ix_j\kappa^d=1/[n(n+1)]$ for $i\ne j$.
The initial velocity is $\sum a_jx_j$, and the smooth geodesic
distance formula gives
\begin{equation}\label{eq:explicit-action}
 \begin{split}
 d_B(\pi^*v_0,\pi^*v_1)^2
      &=\frac{2\bigl[(\sum_ja_j)^2+\sum_ja_j^2\bigr]}{n+1},\\
 \mathcal S_X(\pi^*v)=\mathcal A_B(\pi^*v)
      &=\frac{(\sum_ja_j)^2+\sum_ja_j^2}{n+1}.
 \end{split}
\end{equation}
For $n=3$ and $(a_0,a_1,a_2)=(1,-1,0)$, the squared distance is one
and the action is $1/2$. These paths are spatially nonconstant on the
actual nonproduct Hopf manifold.

\section{An explicit singular finite energy family}
Let $o=[1:0:\cdots:0]\in\PP^d$ and, for $0<a<1$, set canonically
\[
 \psi=\log\frac{\sum_{j=1}^d|z_j|^2}{\sum_{j=0}^d|z_j|^2},
 \qquad v_a=1-(1-\psi)^a,\qquad v_a(o)=-\infty.
\]
In the affine chart write $s=\log|z'|^2$,
$L=\log(1+e^s)$, $P=e^s/(1+e^s)$ and $T=1+L-s\geq1$.
The local potential $\Phi=L+1-T^a$ has slope
\[
 A=P+a(1-P)T^{a-1},\qquad
 A'=P'(1-aT^{a-1})+a(1-a)(1-P)^2T^{a-2}>0.
\]
Thus $v_a$ is $\kappa$-psh off $o$ and smooth across the hyperplane
at infinity. The bounded $\kappa$-psh truncations $\max(v_a,-j)$
decrease to its genuine canonical representative. Nontriviality and
smooth-volume integrability follow from $|v_a|=O((1-s)^a)$ near the
pole and radial volume $O(e^{ds})\,ds$.

The actual radial flux is $Q(s)=A(s)^d$. The sphere coefficient is one:
$d^c\log|z'|^2$ has integral one on its Hopf circle, while its Hessian
is mass-one Fubini--Study on $\PP^{d-1}$. Bounded Stokes on annuli
therefore gives mass $Q(s_2)-Q(s_1)$. One has $Q(-\infty)=0$ and
$Q(+\infty)=1$. Finite-level locality and exhaustion by these annuli
give the actual non-pluripolar mass one, with no pole atom. Smoothness
at infinity gives no mass on that hyperplane.

As $s\to-\infty$, $T\sim1-s$ and
\[
 |v_a|^2\sim T^{2a},\qquad
 Q'(s)\sim d a^d(1-a)T^{-d(1-a)-1}.
\]
Consequently
\begin{equation}\label{eq:energy-cutoff}
 \int_Y|v_a|^2\MA_\kappa(v_a)<\infty
       \quad\Longleftrightarrow\quad a<\frac{d}{d+2}.
\end{equation}
At equality the divergence is logarithmic. Every other compact region
has smooth potential and measure. By \eqref{eq:fiber-integration},
$\pi^*v_a$ has actual mass $2n$ and lies in $\BB_2$ precisely in
\eqref{eq:energy-cutoff}. Its pole set is the entire elliptic fiber over $o$.
For $0<a<b<d/(d+2)$, the two basic potentials have finite positive
$d_B$ distance, but $T^b-T^a\to\infty$ rules out mutual ordinary
comparison. Both have ordinary model projection zero, as follows below.

\section{The ordinary model projection and test curves}
For finite $C$, the full obstacle envelope with obstacle
$\pi^*\min(v+C,0)$ is invariant under $E$. The same descent proof gives
\[
 P_h\bigl(\pi^*\min(v+C,0)\bigr)
       =\pi^*P_\kappa\bigl(\min(v+C,0)\bigr).
\]
Taking the increasing limit and its canonical usc representative yields
$P_h[\pi^*v]=\pi^*P_\kappa[v]$. In particular base ordinary models
embed. For $v\in\mathcal E^2$, the base full-mass envelope theorem
\cite[Theorem 4.3]{Darvas} gives $P_\kappa[v](0)=0$. Hence the
ordinary model projection of every point of $\BB_2$ is zero; the
metric in \eqref{eq:distance} retains potential data that this projection
forgets. Lower and upper volumes of the entire ordinary type include
nonbasic potentials and are not computed by the basic mass formula.

Ordinary-equivalent endpoints have their common type throughout the
Perron segment, by the affine maximum barriers and the chord bound.
No unrestricted stretching-envelope equivalence is inferred from a model
projection. For a base model test curve, pullback preserves concavity,
monotonicity and ordinary-model normalization. It commutes with the forward
usc supremum Legendre transform and the backward pointwise infimum.
If the associated base Perron ray has every finite slice in $\mathcal E^2$,
its pullback has every finite segment a $d_B$ geodesic and metric-action
minimizer. The finite-energy condition is part of this statement.

This construction concerns the complete topology of canonical basic
finite energy potentials on a genuinely nonclosed manifold. A separating
complete metric on all ordinary models, with modification and varying
background compatibility, remains a further problem.

\begin{thebibliography}{9}
\bibitem{BT} E. Bedford and B. A. Taylor,
\emph{A new capacity for plurisubharmonic functions},
Acta Math. \textbf{149} (1982), 1--40.
\bibitem{BTfine} E. Bedford and B. A. Taylor,
\emph{Fine topology, Shilov boundary and $(dd^c)^n$},
J. Funct. Anal. \textbf{72} (1987), 225--251.
\bibitem{BGL} S. Boucksom, V. Guedj and C. H. Lu,
\emph{Volumes of Bott--Chern classes}, arXiv:2406.01090.
\bibitem{Darvas} T. Darvas,
\emph{The Mabuchi completion of the space of K\"ahler potentials},
Amer. J. Math. \textbf{139} (2017), 1275--1313;
arXiv:1401.7318v3.
\end{thebibliography}
\end{document}
